%% presentation.tex
%%
%% Presentation of the course ``System Integration'' of the Official Master on Libre Software (URJC)
%% http://master.libresoft.es
%%

% Compilar a .pdf con LaTeX (pdflatex)
% Es necesario instalar Beamer 
%

\documentclass{beamer}
\usepackage{moreverb} 
\usepackage{listings}
\usepackage{mflogo}
% imprimir
% \documentclass[handout]{beamer} 
% \usepackage{pgfpages}
% \pgfpagesuselayout{4 on 1}[a4paper,landscape,border shrink=5mm]

\mode<presentation> {
  \usetheme{Warsaw}
  \setbeamercovered{opaque}
}

\usebackgroundtemplate{\includegraphics[width=\paperwidth]{format/libresoft-bg.png}}
% \usepackage[spanish]{babel}
\usepackage[utf8]{inputenc}
\usepackage{graphics}
\usepackage{amssymb} % Simbolos matematicos

%% Metadatos del PDF.
\hypersetup{  
  pdftitle={System Integration},
  pdfauthor={Miguel Vidal},
  pdfcreator={GSyC/Libresoft},
  pdfproducer=PDFLaTeX,
  pdfsubject={Master on Free Software},
}
%%

\defbeamertemplate*{footline}{shadow theme}
{%
  \leavevmode%
  \hbox{\begin{beamercolorbox}[wd=.5\paperwidth,ht=2.5ex,dp=1.125ex,leftskip=.1cm plus1fil,rightskip=.2cm]{author in head/foot}%
\includegraphics[scale=0.40]{format/cc-by-80x15.png} \hspace{0.05cm}
	Miguel Vidal / Jose Castro
  \end{beamercolorbox}%
  \begin{beamercolorbox}[wd=.5\paperwidth,ht=2.5ex,dp=1.125ex,leftskip=.3cm,rightskip=.3cm plus1fil]{title in head/foot}%
    \usebeamerfont{title in head/foot}\insertshorttitle%
    \hspace{3.25cm} \usebeamerfont{author in head/foot}\insertframenumber\,/\,\inserttotalframenumber\hfill
  \end{beamercolorbox}}%
  \vskip0pt%
}


\begin{document}

\title{System Integration}
\subtitle{Master on Libre Software (URJC)}
% \institute{\texttt{http://gsyc.urjc.es/\~{}mvidal} \\ Twitter: \texttt{@mvidallopez}}
\author{Miguel Vidal, Jose Castro} 
\date{\footnotesize{Fuenlabrada (Madrid) \\ March-April, 2012}}
\author{Miguel Vidal \hspace{1cm} Jose Castro \\
\hspace{0.5mm} {\tiny Twitter: @mvidallopez \hspace{1.1cm}Twitter: @jfcastroluis}
}


\frame{
\maketitle
\begin{center}
\includegraphics[width=6cm]{format/gsyc-urjc}
\end{center}
}

%% License slide
\begin{frame}
  \vspace{2cm}
  \begin{flushright}
    {\small \copyright{} 2010-2012 Miguel Vidal, Jose Castro} \\
%    \vspace{0.25cm}
    \medskip
    {\scriptsize This work is licensed under \\ a Creative Commons Attribution 3.0 License}
%    \vspace{0.10cm}
  \end{flushright}
  \begin{flushright}
    \href{http://creativecommons.org/licenses/by/3.0/es}{\includegraphics[width=2cm]{format/cc-by.png}} \\
    {\tiny \url{http://creativecommons.org/licenses/by/3.0}}
  \end{flushright}
\end{frame}%%

\usebackgroundtemplate{}

\AtBeginSection[]
{
\begin{frame}<beamer>
\begin{center}
{\huge \insertsection}
\end{center}
\end{frame}
}


\AtBeginSubsection[]
{
  \begin{frame}<beamer>{Table of Contents}
    \tableofcontents[currentsection,currentsubsection]
  \end{frame}
}

%%

\normalsize

%%---------------------------------------------------------------------
%%---------------------------------------------------------------------

\begin{frame}
  \frametitle{Course Contents}

%  \begin{itemize}[<+->]
  \begin{itemize}
    \item \alert{Session 0: Presentation of the Course}
    \item Session 1: Introduction to free software on Servers
    \item Session 2: Storage as a Service
    \item Session 3: The Art of Virtualization
    \item Session 4: From Virtualization to Cloud Computing
    \item Session 5: High Availability Clusters
  \end{itemize}

\end{frame}


\section{Presentation of the course}

%%---------------------------------------------------------------

\begin{frame}
\frametitle{Administrative data}

\begin{itemize}
\item Second semester, 3 ECTS credits
\item Lecturers:
  \begin{itemize}
  \item \alert{Miguel Vidal} (\texttt{mvidal at gsyc.urjc.es})
  \item \alert{Jose Castro} (\texttt{jfcastro at libresoft.es})
  \end{itemize}
\item Schedule: \alert{Fridays, 17:00-21:00}
  \begin{itemize}
  \item Classroom 227, Departamental I, campus at Fuenlabrada
%  \item Lab 7, Laboratorios II, campus at Fuenlabrada
  \end{itemize}
\item Moodle course: \\
  {\footnotesize\url{http://gsyc.escet.urjc.es/moodle/course/view.php?id=130}}
\end{itemize}
\end{frame}


%%---------------------------------------------------------------

\begin{frame}
\frametitle{Free Software basis}

\pause
\begin{itemize}
    \item Development
\pause
    \item Legal issues
\pause
    \item Economics / Business Models
\pause
    \item Community
\pause
    \item \alert{Systems!} 
\end{itemize}

\end{frame}


%%---------------------------------------------------------------

\begin{frame}
\frametitle{Systems Integration?}

\begin{itemize}
   \item Linking together hardware/software subsystems into a whole (one system or ``\alert{platform}'').
   \item Broad range of techniques: computer networking, enterprise application integration, system administration or programming.
   \item System integration engineer: Broad range of skills. Challenging problems.
   \item Systems from a professional standpoint. 
   \item In short: ``\alert{pulls it all together.}'' 
\end{itemize}

\end{frame}



%%---------------------------------------------------------------

\begin{frame}
\frametitle{Main topics addressed}

\begin{itemize}
   \item Overview about some advanced tasks that System Administrators and Network Engineers are typically expected to perform with free software. 
   \item We will learn from design issues to integrate and deployment some complex architectures.
   \item Focus in current SA skills: \alert{Virtualization} (servers, storage, networks) and \alert{clustering}.   
\end{itemize}

\end{frame}


%%---------------------------------------------------------------

\begin{frame}
\frametitle{Sessions}

\begin{enumerate}
\item \alert{March 23rd:} Introduction to free software on Servers
\item \alert{March 30th:} Storage as a Service
\item \alert{April 13rd:} The Art of Virtualization
\item \alert{April 20th:} From Virtualization to Cloud Computing
\item \alert{April 27th:} High Availability Clusters
\end{enumerate}

\end{frame}


%%---------------------------------------------------------------

\begin{frame}
\frametitle{Evaluation criteria}

\begin{itemize}
\item \alert{\textsc{Final Practice}} will represent \alert{100\% of final evaluation}.
\end{itemize}
\end{frame}

%%---------------------------------------------------------------

\begin{frame}
\frametitle{Grades}

\begin{itemize}
\item \textbf{0--4.9: Fail} (``Suspenso'', SS)
\item \textbf{5--6.9: Pass} (``Aprobado'', AP)
\item \textbf{7--8.9: Good} (``Notable'', NT)
\item \textbf{9--10: Excellent} (``Sobresaliente'', SB)
\end{itemize}

\end{frame}

%%---------------------------------------------------------------

\begin{frame}
\frametitle{Final Practice}

\vspace{-1cm}

\begin{center}
\includegraphics[width=6.5cm]{../practica-2012.pdf}
\end{center}

\end{frame}

%%---------------------------------------------------------------

\begin{frame}
\frametitle{Final Practice}

\begin{itemize}
\item  \alert{+1 point:} deploying VMs with two differents free Operating Systems (instead of Linux).
\item  \alert{+1 point:} VMs will have to mount LVM filesystems. So different volumes (LVMs) for different VM mountpoints:
	\begin{itemize}
	\item \texttt{/var/lib/www} volume mounted in \texttt{ws01}
	\item \texttt{/var/lib/mysql} volume mounted in \texttt{db01}
	\end{itemize}
\item  \alert{+1 point:} Doing Moodle exercises and being proactive in classroom.
\item  \alert{+0.5 point:} a web server different to Apache for \texttt{ws01}, such as nginx or Cherokee.
\item  \alert{+0.5 point:} \textsc{https} support in web server (\texttt{ws01}). You'll have to open (and forward) 443 port in your firewall host to \texttt{ws01}.
\end{itemize}
\end{frame}

%%---------------------------------------------------------------

\begin{frame}
\frametitle{Submission deadline}

\begin{itemize}
\item For May evaluation: \alert{May 12th 2012}. 
\item For June/July evaluation: the date limit will be the \alert{29th of June}.
\end{itemize}

\end{frame}


%%---------------------------------------------------------------

\begin{frame}
\frametitle{Some general references}

\begin{itemize}
\item Course Website: \\
  {\footnotesize \url{http://gsyc.escet.urjc.es/moodle/course/view.php?id=130}}
\item \textsc{Evi Nemeth et al.}, \textit{Unix and Linux System Administration Handbook} (4th Edition), Prentice Hall, July 2010.
\item \textsc{Aeleen Frisch}, \textit{Essential System Administration} (3rd Edition), O'Reilly, August 2002.
\item \textsc{Thomas A. Limoncelli} \textit{The Practice of System and Network Administration} (2nd Edition) Addison-Wesley, 2007. 
\item \textsc{Thomas A. Limoncelli}, \textit{Time Management for System Administrators} (3rd Edition), O'Reilly, August 2005.

\end{itemize}

\end{frame}

\frame{
\maketitle
\begin{center}
\includegraphics[width=6cm]{format/gsyc-urjc}
\end{center}
}

\end{document}
